% ============================================================================
% §2.1 Phương pháp — IRT, LSEM, CT-DSEM, Kalman, ứng dụng. Văn phong tiếng Việt
% học thuật, giảm code-switching không cần thiết.
% ============================================================================

\subsection{Phương pháp}
\label{sec:methods}

Khung phương pháp gồm hai tầng (Hình~\ref{fig:pipeline}): tầng \emph{đo lường} sử dụng IRT để biến dữ liệu nhị phân thành ước lượng năng lực; tầng \emph{cấu trúc} sử dụng LSEM để phân tích quỹ đạo của năng lực theo thời gian. Bộ lọc Kalman cục bộ đóng vai trò cầu nối thời gian thực giữa hai tầng, đồng thời tận dụng sai số đo \(\seit\) làm phương sai quan sát biến thiên theo thời gian.

\subsubsection{Tầng đo lường: IRT 1PL và chấm điểm EAP}

Mô hình Rasch một tham số (1PL) xác định xác suất trả lời đúng câu \(j\) của học sinh \(i\):
\begin{equation}
P(Y_{ij}=1 \mid \theta_i, b_j) = \sigma(\theta_i - b_j), \quad \sigma(x) = (1+e^{-x})^{-1}.
\label{eq:1pl}
\end{equation}
Tham số độ khó \(b_j\) được ước lượng bằng phương pháp Maximum Marginal Likelihood thông qua gói \texttt{mirt} \cite{chalmers2012}. Năng lực \(\hat{\theta}_{it}\) tại mỗi đợt đánh giá được chấm điểm bằng phương pháp Expected A Posteriori (EAP) với phân phối tiên nghiệm \(\Norm(0,1)\), kèm sai số chuẩn \(\seit = \sqrt{\Var(\theta\mid\mathbf{y}_{it})}\). Mô hình 1PL được lựa chọn thay cho 2PL/3PL do trên bộ dữ liệu hiện có, 2--6\% câu hỏi cho tham số discrimination thoái hoá khi ước lượng 2PL; trong khi đó, 1PL ổn định trên toàn ngân hàng câu hỏi \cite{lord1980}.

\subsubsection{Tầng cấu trúc: bốn biến thể LSEM}

\textit{Mô hình tăng trưởng tiềm ẩn (LGCM)} \cite{mcardle2009} mô hình hoá đồng thời điểm xuất phát và tốc độ tăng trưởng của cá nhân:
\begin{equation}
\hat{\theta}_{it} = \alpha_i + s_i\,t + \varepsilon_{it}, \quad (\alpha_i, s_i)^\top \sim \Norm\big((\alpha, s)^\top,\ \Sigma\big).
\label{eq:lgcm}
\end{equation}
\textit{Mô hình điểm thay đổi tiềm ẩn (LCSM)} bổ sung hệ số coupling \(\beta\) giữa các thời điểm liên tiếp:
\begin{equation}
\Delta \theta_{it} = \alpha + \beta\,\theta_{i,t-1} + \zeta_{it}.
\label{eq:lcsm}
\end{equation}
\(\beta < 0\) thể hiện cơ chế đuổi kịp (học sinh yếu tăng nhanh hơn học sinh khá); \(\beta \approx 0\) thể hiện trạng thái ổn định. \textit{Mô hình tăng trưởng đa cấp (lmer)} sử dụng phương pháp REML thay vì FIML và đóng vai trò \emph{bằng chứng hội tụ}: kết quả của hai phương pháp ước lượng phải đồng dấu và xấp xỉ về độ lớn để khẳng định phát hiện không phụ thuộc lựa chọn thuật toán. \textit{Mô hình CT-DSEM} (Ornstein--Uhlenbeck \cite{driver2017,asparouhov2018}) được áp dụng cho khối~12 --- nơi khoảng cách thời gian \(\Delta t\) giữa các đợt đánh giá rất không đều:
\begin{equation}
d\theta_{it} = (\mu - \varphi\,\theta_{it})\,dt + \sigma\,dW_t,
\label{eq:ctdsem}
\end{equation}
với điểm cân bằng dài hạn \(\theta_{\mathrm{eq}} = -\mu/\varphi\) và bán chu kỳ \(\ln 2 / |\varphi|\) đặc trưng cho tốc độ tiến về cân bằng. Mô hình RI-CLPM \cite{hamaker2015} không được sử dụng do dữ liệu hiện tại chỉ có một biến tiềm ẩn duy nhất (năng lực Toán \(\theta\)); trong trường hợp đơn biến, RI-CLPM trở nên thừa tham số và quy về một mô hình AR(1) ngẫu nhiên, không cung cấp thêm thông tin so với CT-DSEM. Khoảng tin cậy 95\% cho các hệ số LGCM/LCSM được tính bằng \emph{bias-corrected bootstrap} với 1{.}000 mẫu (lavaan); khoảng tin cậy cho hệ số đa cấp tính theo Wald từ ma trận thông tin Fisher (lme4); khoảng tin cậy cho CT-DSEM tính theo Hessian xấp xỉ tại điểm MAP (ctsem).

\subsubsection{Bộ lọc Kalman cục bộ với phương sai quan sát biến thiên}

Đối với mỗi học sinh \(i\), mô hình không gian-trạng thái tuyến tính được thiết lập như sau: phương trình trạng thái \(\theta_{it} = \theta_{i,t-1} + u_{it}\) với nhiễu trạng thái \(u_{it} \sim \Norm(0,\, q_g)\), và phương trình quan sát \(\hat{\theta}_{it} = \theta_{it} + \varepsilon_{it}\) với \(\Var(\varepsilon_{it}) = \seit^{2}\). Điểm mới so với cách áp dụng Kalman thông thường là phương sai quan sát \(H_t = \seit^2\) được lấy trực tiếp từ Phương trình~\eqref{eq:1pl} qua thuật toán EAP và biến thiên theo số câu hỏi mà học sinh làm trong mỗi đợt. Phương sai trạng thái \(q_g\) được ước lượng chung ở cấp khối bằng phương pháp moment matching: \(\hat q_g = \overline{\Delta\hat{\theta}^2} - \overline{\seit^2 + \mathrm{SE}_{i,t-1}^2}\). Thuật toán làm trơn được thực hiện qua gói \texttt{KFAS::KFS} \cite{helske2017}, ổn định trên cả chuỗi quan sát ngắn (\(T \ge 2\)). Mở rộng sang quan sát phi tuyến (1PL/2PL) thông qua Suy diễn Biến phân được phác thảo trong tài liệu thiết kế kèm theo nhưng chưa triển khai mã chạy thực tế.

\subsubsection{Hai ứng dụng kế tiếp}

\textit{Gợi ý câu hỏi tối đa hoá thông tin (B7b).} Trong mô hình 1PL, hàm thông tin Fisher \(I(\theta,b) = \sigma(\theta-b)(1-\sigma(\theta-b))\) đạt cực đại tại \(b = \theta\). Bộ gợi ý chọn \(k\) câu có \(|b - \hat{\theta}_{i,T}| < \delta\) (mặc định \(\delta = 0{,}5\)) từ ngân hàng câu hỏi của khối tương ứng.

\textit{Cảnh báo sớm học sinh có dấu hiệu giảm năng lực (B7a).} Một học sinh được đánh dấu nếu thoả mãn ít nhất một trong hai tiêu chí: (a) tốc độ tăng trưởng \(s_i < Q_{10}(s)\) --- thuộc nhóm 10\% thấp nhất trong khối; (b) \(\min_t \Delta\theta_{i,t}^{\,\mathrm{smooth}} < -1\) logit --- xuất hiện một bước giảm đột ngột vượt mức 2,5 lần sai số chuẩn sau làm trơn. Tốc độ tăng trưởng cho tiêu chí (a) được trích từ thành phần ngẫu nhiên (\texttt{ranef}) của mô hình lmer, không phải hồi quy OLS từng cá nhân; tín hiệu giảm dùng cho tiêu chí (b) lấy từ quỹ đạo đã làm trơn bằng Kalman thay vì \(\hat\theta\) thô để hạn chế dương tính giả.
